% Motivacija

Графови представљају једну од основних структура података у рачунарству. Једноставност ове апстрактне структуре омогућава велику изражајност и погодност за моделовање различитих система, од различитих мрежа и мапа до односа међу објекатима. Међутим, у практичној имплементацији граф се не користи директно као апстрактан појам, већ се представља конкретним структурама података, као што су листе суседства, матрице суседства или многи компресовани формати за ретке графове \cite{vandergrinten_et_al:LIPIcs.SEA.2022.11}. Избор ових репрезентација значајно утиче на ефикасност алгоритама, како у погледу временске, тако и просторне сложености. Већина стандардних алгоритама и структура за рад са графовима полази од претпоставке да је граф статичан, односно да се скуп чворова и грана не мења током извршавања алгоритма. Насупрот томе, многи реални системи подразумевају промене током времена, где се чворови и гране често додају, уклањају или мењају \cite{APPLICATIONS_OF_DYNAMIC_GRAPH}. Овакви системи захтевају проширење досадашњих погледа на графове.

Узмимо на пример комуникацију у рачунарској мрежи, која се природно може представити графом (слика \ref{fig:network}). Рутери и везе између њих могу се појавити и нестати услед различитих фактора као што су кварови, одржавање или промене у мрежној топологији. Иако се структурне промене у мрежи (додавање или уклањање чворова и грана) могу јављати релативно ретко, уколико се тежинама грана моделују динамичке величине, као што је тренутни пропусни опсег везе, вредности ових тежина се мењају континуирано током времена. У таквом моделу, одржавање ажурних информација о стању мреже, укључујући оптималне путање рутирања и друга релевантна својства, представља значајан алгоритамски изазов.

\begin{figure}[ht]
\centering
\begingroup
\shorthandoff{"}
\digraph[scale=0.5]{network}{
    rankdir=LR;
    edge [dir=none, color=black, splines=line];
    A [label=A, shape=box3d, style=filled, fillcolor=lightblue];
    B [label=B, shape=box3d, style=filled, fillcolor=lightblue];
    r1 [label=R1, shape=circle];
    r2 [label=R2, shape=circle];
    r3 [label=R3, shape=circle];
    r4 [label=R4, shape=circle];
    r5 [label=R5, shape=circle];
    r6 [label=R6, shape=circle];

    A -> r1 [label="10 Gbps"];

    r1 -> r2 [label="5 Gbps"];
    r1 -> r3 [label="2 Gbps"];

    r2 -> r4 [label="4 Gbps"];
    r3 -> r4 [label="1 Gbps"];

    r2 -> r5 [label="3 Gbps"];
    r5 -> r6 [label="2 Gbps"];

    r4 -> r6 [label="1 Gbps"];

    r6 -> B [label="4 Gbps"];
}
\endgroup
\caption{Неусмерена мрежа рутера између два рачунара}
\label{fig:network}
\end{figure}

У таквим динамичким окружењима, алгоритми који раде на статичким графовима нису довољни, јер не узимају у обзир промене које се дешавају током времена и могу довести до нетачних резултата. Поновна анализа комплетне мреже након сваке промене постаје рачунски неизводљива, због чега су неопходни алгоритми и структуре података који се могу ефикасно прилагођавати изменама. Ови приступи омогућавају одржавање ажурних информација о мрежи, брзу реакцију на промене, као и ефикасно решавање задатака као што су праћење мреже, откривање аномалија и оптимизација протока података у окружењима са високим степеном динамике \cite{APPLICATIONS_OF_DYNAMIC_GRAPH}. 

Овај рад се бави наведеним изазовима, кроз анализу алгоритама и структура података за динамичке графове, са циљем проналажења ефикасних решења за проблеме који настају у оваквим променљивим окружењима.

Формална дефиниција динамичког графа биће наведена касније, али интуитивно, динамички граф је граф проширен временском димензијом у којој се мења његова структура. Ове промене могу бити узроковане различитим факторима, зависно од конкретне примене, али уопштено, динамички графови су неопходни за моделовање и анализу променљивих система. Увођењем временског аспекта доводи се у питање не само како ефикасно представити и манипулисати графом, већ и како одржавати информације о његовој структури и својствима.


% Motivacija